% !TEX root = ../DoAn.tex
\documentclass[../DoAn.tex]{subfiles}
\begin{document}

% =======================================================================
\section*{Kết luận}
\addcontentsline{toc}{section}{Kết luận}

Đồ án đã hoàn thành mục tiêu xây dựng một hệ thống thương mại điện tử theo kiến trúc
microservices, đi trọn vẹn từ phân tích yêu cầu, thiết kế, hiện thực hóa, đóng gói, triển khai
cho đến kiểm thử trên môi trường gần thực tế. Hệ thống gồm 13 service nghiệp vụ và 3 service hạ
tầng Spring Cloud, xây dựng trên Java~21, Spring Boot~3.5 và Spring Cloud~2025, đáp ứng đầy đủ
sáu nhóm mục tiêu đặt ra ở Chương~\ref{chapter:Introduction}.

Về kiến trúc, hệ thống được phân rã thành các Bounded Context độc lập theo nguyên tắc
Domain-Driven Design, mỗi service sở hữu cơ sở dữ liệu riêng và chỉ giao tiếp qua API hoặc sự
kiện Kafka, với API Gateway làm cổng vào duy nhất, Eureka khám phá dịch vụ và Config Server quản
lý cấu hình tập trung. Các bài toán khó của giao dịch phân tán được giải quyết bằng những mẫu
thiết kế phù hợp: Saga Orchestration điều phối luồng đặt hàng, Transactional Outbox kết hợp
Idempotent Consumer bảo đảm xử lý exactly-once ở tầng ứng dụng, và CQRS-lite với Elasticsearch
cho tìm kiếm toàn văn. Bảo mật được tập trung tại Gateway theo Trusted Subsystem Pattern dựa trên
Keycloak và JWT RS256, kết hợp rate limiter cho các luồng nhạy cảm. Hệ thống cũng tích hợp thanh
toán VNPAY và một cơ chế flash sale chống over-sell đa lớp dựa trên Redis Lua atomic reservation.

Điểm nổi bật nhất về mặt kỹ thuật là cơ chế flash sale đa lớp và bảo đảm exactly-once ở tầng ứng
dụng --- cả hai đều được kiểm chứng bằng số liệu thực nghiệm thay vì chỉ mô tả lý thuyết. Hệ
thống được đóng gói bằng Docker, triển khai production-like trên AWS EC2 với quy trình build/push
image tự động qua GitHub Actions/GHCR, và được trang bị một lớp quan sát hoàn chỉnh gồm
Prometheus/Grafana cho số liệu, Zipkin cho truy vết phân tán và structured logging. Kết quả kiểm
thử trên môi trường này cho thấy hệ thống vận hành ổn định: catalog soak đạt p95 61,68~ms với
220.956 request không lỗi, checkout stress giữ tỷ lệ tạo đơn 100\%, flash-sale spike 500~VU mua
đúng 100/100 slot không có duplicate buyer, và các kịch bản resilience xác nhận khả năng tự phục
hồi qua outbox replay và graceful degradation. Toàn bộ hệ thống được hoàn thiện bằng giao diện
Next.js gồm storefront và Admin Panel, đủ để demo một sản phẩm hoàn chỉnh.

Bên cạnh các kết quả đạt được, đồ án vẫn còn một số hạn chế cần ghi nhận trung thực. Môi trường
triển khai mới ở mức single-host EC2 nên chưa có high availability; quy trình CI/CD mới dừng ở
build/push/deploy cơ bản; frontend phục vụ tốt cho demo nhưng chưa hỗ trợ upload file qua object
storage hay quản lý người dùng đầy đủ qua Keycloak Admin API; event schema đang quản lý bằng
module dùng chung thay vì schema registry; và hệ thống chưa có log aggregation cùng alerting ở
mức production. Những hạn chế này chính là tiền đề cho các hướng phát triển tiếp theo.

% =======================================================================
\section*{Hướng phát triển}
\addcontentsline{toc}{section}{Hướng phát triển}

Hướng phát triển của đồ án tập trung vào ba nhóm chính. \textit{Trước hết là hoàn thiện nền tảng
vận hành:} chuyển từ single-host Docker Compose sang Kubernetes nhiều node để đạt high
availability và tự co giãn theo tải, đồng thời nâng cấp CI/CD với security scan, quản lý secret
chuyên dụng và chiến lược triển khai canary/blue-green, bổ sung log aggregation và alerting cho
khả năng quan sát ở mức production.

\textit{Tiếp theo là mở rộng chức năng nghiệp vụ:} tích hợp object storage (S3 hoặc MinIO) cho
phép upload ảnh thật thay vì nhập URL, kết nối Admin Panel với Keycloak Admin API để quản lý
người dùng hoàn chỉnh, tách một analytics service riêng theo mô hình CQRS, và xây dựng dịch vụ
gợi ý sản phẩm dựa trên lịch sử mua hàng và hành vi duyệt.

\textit{Cuối cùng là nâng cấp nền tảng kỹ thuật cho quy mô lớn hơn:} áp dụng schema registry và
contract testing để quản lý hợp đồng sự kiện giữa các service, triển khai service mesh cho mTLS
và observability ở tầng mạng, và mở rộng sang ứng dụng di động bằng cách tái sử dụng API Gateway
cùng luồng xác thực OAuth~2.0 sẵn có của Keycloak.

\end{document}
